%%%%%%%%%%%%%%%%%%%%%%%%%%%%%%%%%%%%%%%%%%%%%%%%%%%%%%%%%%%%%%%%%%%%%%
% Informe de prensa - Instituto de Economía - UADE
% IIJP: Índice de Inflación de Jubilados y Pensionados (datos hasta septiembre de 2026)
% Plantilla: Informe de prensa del INECO (estilo newsletter, portada y cajas)
%
% Compilar con pdfLaTeX + BibTeX (Overleaf lo hace solo).
% Archivos que necesita el proyecto de Overleaf:
% - figuras/logo_ineco.jpg (el logo de la plantilla)
% - figuras/iijp_01_canasta.png ... figuras/iijp_06_jubilacion_minima.png: output/figuras/ del repo
%   (python scripts/construir.py) o la carpeta figuras/ del ZIP de cualquier notebook.
% - bibliografia.bib (docs/divulgacion/bibliografia.bib del repo)
% Cifras: salen de python scripts/construir.py (ESTADO.md, seccion 2). Si se actualizan los datos,
% revisar todas las cifras del texto y de los cuadros.
%%%%%%%%%%%%%%%%%%%%%%%%%%%%%%%%%%%%%%%%%%%%%%%%%%%%%%%%%%%%%%%%%%%%%%

\documentclass[10pt,a4paper]{article}

% Al principio del preámbulo
\usepackage[a4paper]{geometry}  % pero no definimos aún
%% -------------- Codificación, idioma y fuentes (pdfLaTeX)
\usepackage{caption} % para usar \captionof{table} dentro de minipage
\usepackage{booktabs}
\usepackage{float}
\usepackage[utf8]{inputenc}
\usepackage[T1]{fontenc}
\usepackage[spanish]{babel}        % idioma español
\usepackage{lmodern}               % fuentes T1 robustas
\usepackage[nopatch=footnote]{microtype} % mejoras tipográficas sutiles (sin el parche de notas al pie: choca con la plantilla)
\usepackage{charter}               % fuente tipo Charter para el cuerpo (opcional)
\frenchspacing
\usepackage{csquotes} % mejora las comillas y compatibilidad con babel


\usepackage{tabularx,booktabs}
\usepackage{makecell}

% --- Citado con natbib (estilo apalike)
\usepackage{natbib}
\bibliographystyle{apalike}

%% -------------- Imágenes
\usepackage{graphicx}
\graphicspath{{figuras/}} % para no repetir "figuras/" en cada includegraphics

%% -------------- Colores y enlaces
\usepackage{xcolor}
\definecolor{INECOblue}{HTML}{0055A4}
\definecolor{INECOgray}{HTML}{4E5D6C}
\definecolor{INECOlight}{HTML}{E6EEF7}

\usepackage[hidelinks,hyperfootnotes=false]{hyperref} % notas al pie sin enlace (en columnas el enlace queda roto)
\hypersetup{
  colorlinks=true,
  linkcolor=INECOblue,
  urlcolor=INECOblue,
  citecolor=INECOblue
}
\urlstyle{same}

%% -------------- Diseño de página (estilo newsletter del template original)
% (no usamos geometry para mantener estilo similar al template original)
\setlength\topmargin{-48pt}
\setlength\headheight{0pt}         % sin encabezado
\setlength\headsep{0pt}
\setlength\marginparwidth{-20pt}
\setlength\textwidth{7.0in}
\setlength\textheight{9.5in}
\setlength\oddsidemargin{-30pt}
\setlength\evensidemargin{-30pt}

%% -------------- Paquetes útiles
\usepackage{amssymb,amsmath}       % símbolos matemáticos
\usepackage{multicol}              % columnas múltiples
\usepackage{marvosym}              % iconos (\Letter, \Mundus, \Telefon, ...)
\usepackage{fancyhdr}              % encabezado/pie
\usepackage{titlesec}              % control de títulos
\usepackage[most]{tcolorbox}       % cajas estilizadas
\usepackage{enumitem}              % listas compactas
\usepackage{parskip}
\usepackage{setspace}
\setstretch{1.7}
\usepackage{indentfirst}
\setlength{\parindent}{1.5em}
\usepackage{siunitx}               % números y porcentajes prolijos (opcional)
\sisetup{
  output-decimal-marker = {,},   % la coma para decimales
  group-separator       = {.},   % el punto para miles
  group-minimum-digits  = 4,     % activa separación de miles a partir de 4 dígitos
  group-digits          = integer
}

\usepackage{caption}

\captionsetup{
  font=small,            % achica el texto del caption (incluye "Figura 1")
  labelfont=small,       % achica también la parte de "Figura 1"
  width=1\linewidth    % fuerza a que el caption se divida en varias líneas width=0.7
}

%% -------------- Fecha personalizable y metadatos
% Uso: \Fecha{23 de septiembre de 2025}
\newcommand{\Fecha}[1]{\if\relax\detokenize{#1}\relax \def\FechaText{\today} \else \def\FechaText{#1} \fi}
% Por defecto usa \today (en español)
\Fecha{Octubre 2026} % <-- poner aquí la fecha si querés fijarla manualmente, ej: \Fecha{23 de septiembre de 2025}

\newcommand{\Numero}[1]{\def\NumeroText{#1}}
\Numero{1} % número de edición

\newcommand{\JournalNameText}{Instituto de Economía - UADE}
\newcommand{\JournalSubtitleText}{Informe de prensa}

%% -------------- NUEVO: Título del informe y autor/investigador
% Definiciones con texto por defecto vacío para evitar errores si se omiten.
\def\TituloInformeText{}
\def\AutorInformeText{}
\newcommand{\TituloInforme}[1]{\def\TituloInformeText{#1}}
\newcommand{\AutorInforme}[1]{\def\AutorInformeText{#1}}

\TituloInforme{%
{\setstretch{1} % interlineado normal SOLO en el título
\noindent       % evita sangría
\huge Inflación de los jubilados \\[0.2cm]
{\LARGE\color{INECOgray}Desde noviembre de 2023, la canasta de los jubilados se encareció un 3,1\% más que la del IPC}
\par}
}
% \AutorInforme{...}  % opcional: nombre del investigador o del equipo

%% -------------- Encabezado / Pie (sin encabezado; pie refinado)
\pagestyle{fancy}
\fancyhf{}                          % limpiar cabeceras/pies
\renewcommand{\headrulewidth}{0pt}  % sin línea superior

% Pie con línea tenue, email/enlace a la izquierda y número de página a la derecha
\renewcommand{\footrulewidth}{0.4pt}
\renewcommand{\footrule}{\hbox to\headwidth{\color{INECOlight}\leaders\hrule height \footrulewidth\hfill}}

\fancyfoot[L]{%
  \sffamily\footnotesize\color{INECOgray}%
  \Letter\ \texttt{investigacion@uade.edu.ar}\\[2pt]
  \Mundus\ \href{https://www.uade.edu.ar/sites/investigacion/instituto-de-economia-ineco/}{INECO}
}
\fancyfoot[R]{\sffamily\footnotesize\color{INECOgray} Página \thepage}

%% -------------- Macros y estilo (reglas, separadores, títulos)
\newcommand{\HorRule}[1]{\noindent\rule{\linewidth}{#1}}
\newcommand{\SepRule}{\noindent\begin{center}\rule{250pt}{1pt}\end{center}}

% Título principal (no se usa directamente, lo reemplaza la portada)
\newcommand{\JournalName}[1]{%
    \begin{center}
      {\sffamily\Huge\bfseries #1}
    \end{center}
    \par \normalsize \normalfont}

% Encabezado de edición (opcional, ya no se usa con \CoverHeader)
\newcommand{\JournalIssue}[1]{%
    \begin{flushright}\small \FechaText \end{flushright}\par\normalsize\noindent}

% Estilo de secciones
\titleformat{\section}
  {\sffamily\Large\bfseries\color{INECOblue}}
  {}{0pt}{}
\titlespacing*{\section}{0pt}{6pt}{4pt}

% Noticia / item (usado como estilo de subtítulo dentro de cajas)
\newcommand{\NewsItem}[1]{%
    \sffamily\large #1 \par\normalsize \normalfont \vspace{4pt}}
\newcommand{\NewsAuthor}[1]{%
    \hfill \textsc{#1} \vspace{4pt} \par \normalfont}

% Caja estándar para bloques destacados (como Resumen ejecutivo)
\tcbset{
  boxrule=0pt, frame hidden,
  colback=INECOlight,
  enhanced, sharp corners,
  borderline west={3pt}{0pt}{INECOblue},
  left=8pt, right=8pt, top=8pt, bottom=8pt
}

% Portada con logo grande arriba y encabezado fecha + subtítulo
\newcommand{\CoverHeader}{%
  % Empujar todo hacia arriba, eliminando margen superior
  \vspace*{-2cm}%

  % --- LOGO ---
  \noindent
  \begin{center}
    \includegraphics[width=\linewidth,keepaspectratio]{logo_ineco.jpg}%
  \end{center}

  \vspace{0.1cm}

  % --- Fecha grande a la derecha + subtítulo a la izquierda ---
  \noindent
  \begin{minipage}[t]{0.48\linewidth}
    {\sffamily\bfseries\large\color{INECOgray} Informe de prensa}
  \end{minipage}%
  \hfill
  \begin{minipage}[t]{0.48\linewidth}
    \raggedleft {\sffamily\bfseries\large\color{INECOgray} \FechaText}
  \end{minipage}

  % --- Título del informe y autor ---
  \vspace{0.3cm}
  {\centering
    {\sffamily\bfseries\fontsize{22}{26}\selectfont \TituloInformeText\par}
    \vspace{0.1cm}
    {\sffamily\large \AutorInformeText\par}
  }%
}


% Columnas con separación y regla tenue
\setlength{\columnsep}{18pt}
\setlength{\columnseprule}{0.2pt}
\makeatletter
\def\columnseprulecolor{\color{INECOlight}}
\makeatother

% Listas compactas
\setlist[itemize]{left=0pt,label=•,itemsep=2pt,topsep=2pt}
\setlist[enumerate]{left=0pt,itemsep=2pt,topsep=2pt}

%% -------------- Inicio del documento
\begin{document}

% ---------------------------
% PORTADA
% ---------------------------
\CoverHeader

% ---------------------------
% BARRAS SUPERIORES FINAS (detalle visual)
% ---------------------------
\noindent\HorRule{3pt} \\[-0.75\baselineskip]
\HorRule{1pt}

% ---------------------------
% RESUMEN EJECUTIVO (en caja)
% ---------------------------
\vspace{0.1cm}
\begin{center}
\begin{minipage}[h]{0.85\linewidth}
\begin{tcolorbox}
    {\normalsize
    \NewsItem{\textbf{\Large Resumen ejecutivo}}
    \noindent
    \begin{itemize}
        \item Según el \textbf{Índice de Inflación de Jubilados y Pensionados (IIJP)} de INECO, \textbf{desde noviembre de 2023 los precios de la canasta de los jubilados subieron un 3,1\% más que los del IPC}. Para quienes cobran la mínima, la diferencia llega a 4,7\%.
        \vspace{.25 cm}
        \item En septiembre de 2026 el IIJP aumentó \textbf{1,9\% en el mes} (IPC: 1,9\%) y acumula \textbf{34,6\% en los últimos doce meses} (IPC: 33,4\%).
        \vspace{.25 cm}
        \item La brecha se explica sobre todo porque \textbf{el IPC todavía usa la canasta de consumo de 2004/05}: pondera mucho la ropa, que aumentó menos que el promedio, y poco la electricidad y el gas, que aumentaron más.
        \vspace{.25 cm}
        \item Con el bono congelado en \$70.000 desde marzo de 2024, \textbf{la jubilación mínima con bono compraba en septiembre de 2026 un 9,8\% menos que en noviembre de 2023} según el IPC, y un 13,9\% menos según la inflación de los jubilados que cobran la mínima.
        \vspace{.25 cm}
        \item Si desde abril de 2024 las jubilaciones se hubieran ajustado por el IIJP, \textbf{en septiembre de 2026 la mínima sería un 1,1\% más alta (\$4.698 por mes)}, con un costo fiscal cercano al \textbf{0,02\% del PIB por año}.
    \end{itemize}
    }
\end{tcolorbox}
\end{minipage}
\end{center}

\clearpage
\section*{La inflación de los jubilados}

El Índice de Precios al Consumidor (IPC) del INDEC mide cuánto cambia el costo de una canasta que representa el consumo promedio de todos los hogares \citep{INDEC2026_IPC}. Pero los jubilados destinan una parte mayor de su presupuesto que el resto de los hogares a alimentos, medicamentos y servicios de la vivienda, y una parte menor a educación, ropa o comidas fuera de casa. Además, el IPC oficial todavía usa la estructura de gasto de la Encuesta Nacional de Gastos de los Hogares (ENGHo) de 2004/05\footnote{Las ponderaciones del IPC surgen de la ENGHo 2004/05, actualizadas por precios a diciembre de 2016 \citep{INDEC2019_IPC_metodologia}. Su actualización con la ENGHo 2017/18 todavía no se implementó.}. Como desde 2024 las jubilaciones se actualizan cada mes con la inflación del IPC de dos meses antes \citep{PEN2024_DNU274}, \textbf{medir la inflación que enfrentan los jubilados permite saber si ese ajuste alcanza}.

El \textbf{Índice de Inflación de Jubilados y Pensionados (IIJP)}, elaborado por INECO, valoriza cada mes, con los precios que releva el INDEC, la canasta de consumo de los hogares que viven de una jubilación o pensión. \textbf{En septiembre de 2026 el IIJP aumentó 1,9\%, lo mismo que el IPC (1,9\%). En los últimos doce meses acumula 34,6\%, frente al 33,4\% del IPC}. Para los jubilados que cobran la mínima, la inflación fue de 2,1\% en el mes y de 36,6\% en los últimos doce meses (cuadro \ref{tab:mes}).

\begin{table}[H]
\centering
\renewcommand{\arraystretch}{1.15}
\caption{Inflación de septiembre de 2026 (\%)}
\label{tab:mes}
\begin{tabular}{@{}l c c c@{}}
\toprule
 & \textbf{Mensual} & \textbf{Últimos 12 meses} & \textbf{Acumulada en 2026} \\
\midrule
IPC oficial (INDEC)          & 1,9 & 33,4 & 23,6 \\
IIJP                         & 1,9 & 34,6 & 24,4 \\
IIJP de la jubilación mínima & 2,1 & 36,6 & 26,2 \\
\bottomrule
\end{tabular}\\[0.2cm]
\caption*{Fuente: Instituto de Economía de UADE (INECO) en base a IPC y ENGHo 2017/18 (INDEC).}
\end{table}

\section*{Los jubilados consumen distinto}

La canasta de los jubilados surge de los gastos que registra la ENGHo 2017/18 \citep{INDEC2020_ENGHo}. Los hogares jubilados destinan el 14,5\% de su gasto a vivienda y servicios (9,4\% en la canasta del IPC), el 11,4\% a salud (8,0\%) y el 5,2\% a comunicación (2,8\%). En cambio, \textbf{gastan menos en prendas y calzado (5,2\% contra 9,9\%), en restaurantes y hoteles (4,5\% contra 9,0\%) y en educación (0,8\% contra 2,3\%)}. En alimentos el peso es parecido al del IPC (26,4\% contra 26,9\%), aunque mayor que el del conjunto de los hogares en la misma encuesta (22,7\%). Dentro de salud, más de la mitad del gasto de los jubilados se va en medicamentos, y dentro de vivienda, la mitad en electricidad y gas (figura \ref{fig:canasta}).

\begin{figure}[H]
\centering
\includegraphics[width=0.75\linewidth]{iijp_01_canasta.png}
\caption{Cómo se reparte el gasto de consumo: canasta del IPC oficial y canasta de los hogares jubilados (\% del gasto).}
\label{fig:canasta}
\end{figure}

\section*{Año por año}

El IIJP no siempre sube más que el IPC (cuadro \ref{tab:anual}). Entre 2017 y 2019 y en 2023 y 2024 los precios de la canasta de los jubilados aumentaron más que los del IPC; entre 2020 y 2022, menos. En 2025 prácticamente no hubo diferencia, y en lo que va de 2026 el IIJP vuelve a estar por encima.

\begin{table}[H]
\centering
\renewcommand{\arraystretch}{1.15}
\caption{Inflación anual (diciembre a diciembre, \%)}
\label{tab:anual}
\begin{tabular}{@{}l c c c c@{}}
\toprule
\textbf{Año} & \textbf{IPC oficial} & \textbf{IIJP} &
\shortstack{\textbf{IIJP}\\\textbf{jubilación mínima}} &
\shortstack{\textbf{Diferencia}\\\textbf{IIJP -- IPC}} \\
\midrule
2017 & 24,8  & 26,4  & 27,4  & $+$1,3 \\
2018 & 47,6  & 49,6  & 50,2  & $+$1,4 \\
2019 & 53,8  & 54,6  & 54,8  & $+$0,5 \\
2020 & 36,1  & 33,9  & 34,2  & $-$1,7 \\
2021 & 50,9  & 48,3  & 46,8  & $-$1,8 \\
2022 & 94,8  & 92,1  & 91,2  & $-$1,4 \\
2023 & 211,4 & 213,3 & 218,2 & $+$0,6 \\
2024 & 117,8 & 121,6 & 121,4 & $+$1,8 \\
2025 & 31,5  & 31,4  & 31,2  & $-$0,1 \\
2026 (enero a septiembre) & 23,6 & 24,4 & 26,2 & $+$0,6 \\
\bottomrule
\end{tabular}\\[0.2cm]
\caption*{Diferencia: cuánto más (o menos) se encareció la canasta de los jubilados que la del IPC, en \%. Fuente: Instituto de Economía de UADE (INECO) en base a IPC y ENGHo 2017/18 (INDEC).}
\end{table}

\section*{Una historia en tres etapas}

Para comparar dos índices se usa la \emph{brecha}: cuánto más se encareció una canasta que la otra\footnote{Brecha $= (1 + \text{inflación del IIJP}) / (1 + \text{inflación del IPC}) - 1$.}. Un ejemplo: entre noviembre de 2023 y septiembre de 2026 el IPC subió 344,4\% y el IIJP 358,0\%. Lo que en noviembre de 2023 costaba \$100, en septiembre de 2026 costaba \$444 según el IPC y \$458 según el IIJP: un 3,1\% más.

Si se mira todo el período, desde diciembre de 2016 hasta septiembre de 2026, la canasta de los jubilados se encareció apenas un 1,2\% más que la del IPC. Pero ese número resume \textbf{tres etapas muy distintas} (figura \ref{fig:historia}): entre diciembre de 2016 y diciembre de 2019 la brecha fue de $+$3,2\%; entre diciembre de 2019 y noviembre de 2023, de $-$4,9\%, y desde noviembre de 2023 es de $+$3,1\%.

\begin{figure}[H]
\centering
\includegraphics[width=0.8\linewidth]{iijp_02_historia.png}
\caption{Brecha acumulada del IIJP sobre el IPC oficial desde diciembre de 2016 (\%).}
\label{fig:historia}
\end{figure}

Lo que mueve la brecha son, sobre todo, dos rubros: \textbf{la ropa, que pesa mucho en la canasta del IPC y poco en la de los jubilados, y las tarifas de la vivienda (electricidad, gas, agua), que pesan más para los jubilados}. Cuando las tarifas se congelaron y la ropa aumentó más que el promedio (2020--2023), la canasta de los jubilados se encareció menos que la del IPC. Cuando las tarifas se recompusieron y la ropa aumentó menos que el promedio (2017--2019 y desde 2024), se encareció más.

\section*{Desde noviembre de 2023}

La brecha se abrió rápido, entre diciembre de 2023 y abril de 2024, se mantuvo entre 2\% y 3\% durante el resto de 2024 y 2025, y en 2026 volvió a subir: \textbf{en septiembre llega a 3,1\%} (figura \ref{fig:nov23}). Para los hogares jubilados que cobran la mínima es mayor, de 4,7\%, porque su canasta pesa más en alimentos, electricidad y gas, transporte público y medicamentos.

\begin{figure}[H]
\centering
\includegraphics[width=0.8\linewidth]{iijp_03_desde_nov23.png}
\caption{Cuánto más subieron los precios de la canasta de los jubilados que los del IPC oficial desde noviembre de 2023 (\%).}
\label{fig:nov23}
\end{figure}

¿Significa esto que los jubilados enfrentan más inflación que el resto de los hogares? No exactamente. Si se calcula el mismo índice con la canasta de \emph{todos} los hogares de la ENGHo 2017/18, la diferencia con el IPC oficial es todavía mayor: 3,7\% desde noviembre de 2023. Es decir, \textbf{la mayor parte de la brecha no se debe a que los jubilados consuman distinto, sino a que el IPC oficial usa una canasta desactualizada, que subestima la inflación de todos los hogares}. Los jubilados que cobran la mínima, en cambio, sí enfrentaron más inflación que el jubilado promedio: 1,6\% más desde noviembre de 2023.

\section*{¿Qué rubros explican la diferencia?}

La figura \ref{fig:explica} reparte la brecha desde noviembre de 2023 entre los rubros de la canasta. Los que más suman son:

\begin{itemize}
    \item \textbf{Prendas y calzado}: la canasta del IPC le da casi el doble de peso que la de los jubilados (9,9\% contra 5,2\%), y la ropa aumentó mucho menos que el promedio (164\% contra 344\%).
    \item \textbf{Vivienda y servicios}: pesa más en la canasta de los jubilados (14,5\% contra 9,4\%) y aumentó mucho más que el promedio (661\%), por la recomposición de las tarifas de electricidad y gas.
\end{itemize}

Del otro lado, restan los restaurantes y la educación, que aumentaron más que el promedio y pesan menos en la canasta de los jubilados, y la recreación.

\begin{figure}[H]
\centering
\includegraphics[width=0.75\linewidth]{iijp_04_que_explica.png}
\caption{Aporte de cada rubro a la brecha del IIJP con el IPC oficial, noviembre de 2023 a septiembre de 2026 (puntos porcentuales). Los aportes se calculan contra una réplica del IPC hecha con los mismos índices de precios y suman 3,0 puntos.}
\label{fig:explica}
\end{figure}

\section*{Diferencias por región}

El INDEC calcula el IPC en seis regiones. Si se compara el IIJP de cada región con el IPC de la misma región, \textbf{desde noviembre de 2023 la brecha es positiva en todas}: va de 2,0\% en la región Pampeana y 2,7\% en el Gran Buenos Aires a 5,7\% en el Noroeste y 9,3\% en el Noreste (figura \ref{fig:regiones}). Fuera del Gran Buenos Aires y la región Pampeana la muestra de hogares jubilados es más chica, por lo que esas cifras son menos precisas.

\begin{figure}[H]
\centering
\includegraphics[width=0.75\linewidth]{iijp_05_regiones.png}
\caption{Brecha del IIJP de cada región sobre el IPC de la misma región, noviembre de 2023 a septiembre de 2026 (\%).}
\label{fig:regiones}
\end{figure}

\section*{La jubilación mínima}

La jubilación mínima se actualiza todos los meses por inflación, pero \textbf{el bono previsional está congelado en \$70.000 desde marzo de 2024}. Medido en poder de compra (figura \ref{fig:minima}), el haber mínimo con bono cayó fuerte a comienzos de 2024 (en febrero de ese año compraba un 25\% menos que en noviembre de 2023), se recuperó a mediados de 2024 y desde entonces vuelve a perder, porque el bono no se actualiza. \textbf{En septiembre de 2026 compraba un 9,8\% menos que en noviembre de 2023 según el IPC, y un 13,9\% menos si se usa la inflación de los jubilados que cobran la mínima}.

Sin el bono, el haber mínimo está un 10,3\% por encima de noviembre de 2023 según el IPC, y un 5,3\% por encima según el IIJP de la jubilación mínima. Para tener en septiembre de 2026 el mismo poder de compra que en marzo de 2024, el bono debería haber sido de \$198.367.

\begin{figure}[H]
\centering
\includegraphics[width=0.8\linewidth]{iijp_06_jubilacion_minima.png}
\caption{Poder de compra de la jubilación mínima más el bono (noviembre de 2023 = 100).}
\label{fig:minima}
\end{figure}

\section*{¿Y si las jubilaciones se ajustaran por el IIJP?}

Si desde abril de 2024 los haberes se hubieran actualizado con el IIJP en lugar del IPC, con la misma regla de movilidad, \textbf{en septiembre de 2026 la jubilación mínima sería de \$433.331 en lugar de \$428.633: un 1,1\% más (\$4.698 por mes)}. Con el IIJP de la jubilación mínima la diferencia sería de 2,6\% (\$11.188 por mes).

El costo fiscal sería acotado: 0,42\% del gasto en jubilaciones y pensiones en 2024 (abril a diciembre) y 0,22\% en 2025, \textbf{alrededor del 0,02\% del PIB por año}. En los primeros ocho meses de 2026 sube a 0,63\% del gasto, porque la diferencia entre ambos índices se fue acumulando.

En síntesis, desde noviembre de 2023 la canasta de los jubilados se encareció más que la del IPC, sobre todo porque el IPC sigue usando una canasta de hace dos décadas. \textbf{Con el bono congelado, la jubilación mínima pierde poder de compra}, y ajustarla por el IIJP tendría un costo fiscal acotado.

\clearpage
\section*{Metodología}

El IIJP, elaborado por INECO, mide la variación mensual de los precios de la canasta de consumo de los hogares jubilados de la Argentina. La serie es mensual desde diciembre de 2016 (base 100) y se actualiza con cada publicación del IPC.

\textbf{Población.} Se consideran hogares jubilados aquellos en los que las jubilaciones y pensiones, contributivas y no contributivas, aportan al menos la mitad del ingreso del hogar. En los microdatos de la ENGHo 2017/18 \citep{INDEC2020_ENGHo} hay 5.350 hogares de este tipo, que representan a 2,96 millones de hogares de todo el país. El IIJP de la jubilación mínima usa la canasta de los hogares jubilados en los que nadie cobra más de 1,1 haberes mínimos \citep{MECON2026_haber_minimo}: 2.081 hogares en la muestra, que representan a 1,17 millones de hogares (el 40\% de los hogares jubilados).

\textbf{Canasta.} La encuesta registra el gasto de cada hogar en cada producto y servicio. Ese gasto se agrega con los factores de expansión de la encuesta (cada hogar pesa según cuánto gasta, como en el IPC) y se clasifica según los rubros del IPC.

\textbf{Precios.} Cada componente de la canasta se valoriza con el índice de precios más detallado que publica el INDEC en cada una de las seis regiones del país \citep{INDEC2026_IPC}. La parte de cada rubro que no tiene un índice propio publicado se valoriza con un índice deducido de las ponderaciones oficiales del IPC \citep{INDEC2019_IPC_metodologia}. El 83\% de la canasta de los jubilados se valoriza con índices publicados por el INDEC.

\textbf{Índice.} Como el IPC, el IIJP es un índice de canasta fija: sigue cuánto cuesta comprar cada mes las cantidades que consumían los hogares jubilados en 2017/18, el criterio que recomienda el manual internacional del IPC \citep{IMF2020_CPI_manual}. El índice nacional suma las canastas de las seis regiones, cada una según el gasto de sus hogares jubilados.

\textbf{Haberes y costo fiscal.} El haber mínimo proviene de datos.gob.ar \citep{MECON2026_haber_minimo} y el bono, de los decretos publicados en el Boletín Oficial. La simulación aplica desde abril de 2024 la misma regla de movilidad del DNU 274/2024 \citep{PEN2024_DNU274}, con el IIJP en lugar del IPC. El costo fiscal se calcula sobre el gasto en jubilaciones y pensiones del Sector Público Nacional \citep{MECON2026_IMIG}.

\textbf{Confiabilidad.} El IIJP usa los mismos precios que el INDEC: lo que cambia es la canasta. Con la canasta oficial, el mismo método reproduce el IPC nacional con un error máximo de 0,08\% desde 2016. El índice se recalculó de 20 maneras distintas, cambiando de a una por vez las decisiones metodológicas (cómo se define un hogar jubilado, la fórmula del índice, el nivel de detalle de los precios, entre otras): en todas, la canasta de los jubilados se encareció más que la del IPC desde noviembre de 2023, con una brecha de entre 2,4\% y 6,7\%. En cambio, la diferencia acumulada desde 2016 ($+$1,2\%) va de $-$0,2\% a 4,0\% según la variante, por lo que no permite afirmar que los jubilados hayan tenido más inflación que el promedio en los últimos diez años.

\textbf{Limitaciones.} La canasta es la de 2017/18 y no capta cambios de hábitos posteriores. Usa los precios del IPC, por lo que no registra descuentos propios de los jubilados (como la cobertura del PAMI en medicamentos, la tarifa social del transporte o la segmentación de los subsidios a la energía), y el IPC no releva precios de geriátricos ni de residencias para mayores.

El código, los datos y la metodología completa son públicos y permiten reproducir todos los resultados \citep{INECO2026_IIJP_repo}.

\clearpage
\bibliography{bibliografia}

\end{document}
